% Publication figures from paper/scripts/render_paper_figures.py.
% Captions folded from paper/harness/phase4_viz/FIGURES.md (honesty-aligned).

\begin{figure*}[t]
\centering
\includegraphics[width=0.98\textwidth]{figures/fig1_architecture}
\caption{\textit{DVXR serving path: ingest, encode, route or fuse, calibrate, and explain.}
Schematic of the locked parent-DVXR path (not a performance plot). Each modality is
encoded by a specialist and either routed to the winning single-modality head or
offered to the fusion library. The language model narrates a frozen number and never
computes risk.
\textit{Note:} Nature \emph{Sci.\ Rep.} attention fusion is the same-person EEG+ECG
stress success case and is not a glucose claim.}
\label{fig:architecture}
\end{figure*}

\begin{figure*}[t]
\centering
\includegraphics[width=0.98\textwidth]{figures/fig2_fusion_vs_single}
\caption{\textit{Learned cross-modal fusion versus the best single-modality specialist
on six subject-held-out mental-health and BCI tasks.}
Source: \texttt{presentation/data/goal3\_ablation/comparative\_performance.csv}.
Panel A: AUROC. Panel B: signed $\Delta$ (fusion $-$ best single). Fusion loses or
ties on every row shown (Holm $p{=}1.0$). The depression bars are $0.918$ vs.\ $0.795$,
not the LaBraM $0.961$ identity-leakage-confounded upper bound.
\textit{Note:} The only complementary-signal win in that CSV is same-subject glucose
(CGMacros RMSE@30 $13.33\rightarrow 12.99$\,mg/dL), shown in Fig.~\ref{fig:glucose}.}
\label{fig:fusion}
\end{figure*}

\begin{figure*}[t]
\centering
\includegraphics[width=0.98\textwidth]{figures/fig3_glucose_ladder}
\caption{\textit{Patient-disjoint 30-minute glucose model ladder on CGMacros
(lower is better).}
Source: \texttt{presentation/data/goal2\_fusion/glucose\_model\_ladder.csv}
(\texttt{horizon\_minutes=30}). Gradient boosting (RMSE $12.48$\,mg/dL) edges
NeuroGlycemicNet ($12.99$). Persistence is the required baseline ($17.40$).
\textit{Note:} Reported as measured; the deep net is not dressed as a win. Meal-event
fusion is a separate same-subject ablation, not a depth claim.}
\label{fig:glucose}
\end{figure*}

\begin{figure*}[t]
\centering
\includegraphics[width=0.98\textwidth]{figures/fig4_nature_contrast}
\caption{\textit{When attention-weighted fusion is licensed, and when it is not.}
Kumar et al.\ report a same-person, synchronized EEG+ECG stress result
(\emph{Sci.\ Rep.} 16:15188, 2026). That inductive bias is not licensed for unmatched
EEG$\times$CGM$\times$EHR triples; this paper refuses that join.
\textit{Note:} Nature fusion is a stress-method contrast, not a glucose method.
Diabetes-attention-fusion glycemic fused macro-F1 is excluded (stale).}
\label{fig:nature}
\end{figure*}

\begin{figure*}[t]
\centering
\includegraphics[width=0.98\textwidth]{figures/fig5_headline_specialists}
\caption{\textit{Headline specialists under subject-held-out evaluation
(window-level AUROC with intervals from Table~\ref{tab:headline}).}
Depression $0.961$ is an identity-leakage-confounded upper bound, not a validated
biomarker (subject identity decodable at $88.8\%$, $52\times$ chance).
\textit{Note:} Research-grade screening, not a diagnosis.}
\label{fig:headline}
\end{figure*}

\begin{figure*}[t]
\centering
\includegraphics[width=0.98\textwidth]{figures/fig6_glucose_ablation}
\caption{\textit{CGMacros missing-modality ablation at 30 min (patient-held-out).}
Source: \texttt{presentation/data/goal3\_ablation/leave\_one\_modality\_out\_cgmacros.csv}.
\textit{Note:} Meal events added to CGM are the only complementary-signal win;
removing CGM collapses RMSE and forces abstention. Values are those in the CSV.}
\label{fig:glucose-ablation}
\end{figure*}
